\section{The exact diagonal Mellin ledger}
\label{sec:tpc52-diagonal-mellin}

We now evaluate the canonical envelope on the distinguished physical
direction.  The result is an exact coordinatewise identity.  It uses
neither orthogonality between different Mellin parameters nor a lower
frame for arbitrary signed coefficient densities.

\subsection{One direct cell and \texorpdfstring{$q$}{q} smooth factors}

Fix an integer \(q\geq1\), and let
\(0\ne W_\nu\in C_c^\infty((0,\infty))\) for
\(1\leq\nu\leq q\).  We use the real-axis Mellin convention
\begin{align}
 \widetilde W_\nu(t)
 &:=\int_0^\infty W_\nu(x)x^{-it}\,\frac{dx}{x},
 \label{eq:tpc52-Mellin-transform}\\
 W_\nu(x)
 &=\frac1{2\pi}\int_\R
     \widetilde W_\nu(t)x^{it}\,dt.
 \label{eq:tpc52-Mellin-inversion}
\end{align}
Indeed, \(s\mapsto W_\nu(\e^s)\) is a compactly supported smooth
function on \(\R\), so \(\widetilde W_\nu\) is Schwartz and
\eqref{eq:tpc52-Mellin-inversion} is ordinary Fourier inversion in
the logarithmic coordinate.

Let \(w_B:\R^q\to[1,\infty)\) be measurable and of at most polynomial
growth; more precisely, for a fixed \(B\geq0\), assume
\begin{equation}
 w_B(\bm t)\leq C_B(1+|\bm t|)^B.
 \label{eq:tpc52-Mellin-profile-weight}
\end{equation}
Define the positive product measure
\begin{equation}
 d\mu_{\mathbf W}(\bm t)
 :=(2\pi)^{-q}\prod_{\nu=1}^q
       |\widetilde W_\nu(t_\nu)|\,d\bm t
 \label{eq:tpc52-Mellin-variation-measure}
\end{equation}
and the polar phase
\begin{equation}
 \varsigma_{\mathbf W}(\bm t)
 :=\prod_{\nu=1}^q
   \begin{cases}
    \widetilde W_\nu(t_\nu)/
       |\widetilde W_\nu(t_\nu)|,
       &\widetilde W_\nu(t_\nu)\ne0,\\
    0,&\widetilde W_\nu(t_\nu)=0.
   \end{cases}
 \label{eq:tpc52-Mellin-polar-phase}
\end{equation}
Its value on the \(\mu_{\mathbf W}\)-null zero set is immaterial.  The
fixed weighted Wiener--Mellin mass is
\begin{equation}
 A_B(\mathbf W)
 :=\int_{\R^q}w_B(\bm t)\,d\mu_{\mathbf W}(\bm t)
 =\frac1{(2\pi)^q}\int_{\R^q}
   w_B(\bm t)\prod_{\nu=1}^q
   |\widetilde W_\nu(t_\nu)|\,d\bm t.
 \label{eq:tpc52-weighted-Mellin-mass}
\end{equation}
Schwartz decay and \eqref{eq:tpc52-Mellin-profile-weight} give
\begin{equation}
 0<A_B(\mathbf W)<\infty.
 \label{eq:tpc52-Mellin-mass-finite}
\end{equation}

Let \(\cI\) be a finite nonempty coordinate set.  For
\(i\in\cI\), fix \(c_i\in\C\) and positive samples
\(x_{i,1},\ldots,x_{i,q}>0\).  Put
\begin{equation}
 A_{\bm t}(c)_i
 :=\varsigma_{\mathbf W}(\bm t)c_i
   \prod_{\nu=1}^q x_{i,\nu}^{it_\nu}
 \in\ell^2(\cI).
 \label{eq:tpc52-diagonal-Mellin-layer}
\end{equation}
Every sample phase in \eqref{eq:tpc52-diagonal-Mellin-layer} has unit
modulus.

\begin{theorem}[Exact \(q\)-factor diagonal Mellin identity]
\label{thm:tpc52-diagonal-Mellin}
\label{thm:tpc52-diagonal-ledger}
For \(c=(c_i)_{i\in\cI}\in\ell^2(\cI)\), define
\begin{align}
 \cS_{\mathbf W}c
 &:=\int_{\R^q}A_{\bm t}(c)\,
       d\mu_{\mathbf W}(\bm t),
 \label{eq:tpc52-diagonal-Mellin-synthesis}\\
 \cE_B(\mathbf W;c)
 &:=\int_{\R^q}w_B(\bm t)
       \|A_{\bm t}(c)\|_{\ell^2(\cI)}
       \,d\mu_{\mathbf W}(\bm t).
 \label{eq:tpc52-diagonal-Mellin-envelope}
\end{align}
Then the Bochner integral is absolutely convergent and
\begin{align}
 (\cS_{\mathbf W}c)_i
 &=c_i\prod_{\nu=1}^qW_\nu(x_{i,\nu}),
 \label{eq:tpc52-diagonal-output}\\
 \|\cS_{\mathbf W}c\|_2^2
 &=\sum_{i\in\cI}|c_i|^2
   \prod_{\nu=1}^q|W_\nu(x_{i,\nu})|^2,
 \label{eq:tpc52-diagonal-output-energy}\\
 \cE_B(\mathbf W;c)
 &=A_B(\mathbf W)\|c\|_2.
 \label{eq:tpc52-exact-envelope-equality}
\end{align}

For \(c\ne0\), define the physical profile-overlap ratio
\begin{equation}
 \Theta(\mathbf W;c)
 :=\frac{
   \displaystyle\sum_{i\in\cI}|c_i|^2
    \prod_{\nu=1}^q|W_\nu(x_{i,\nu})|^2}
   {\displaystyle\sum_{i\in\cI}|c_i|^2}.
 \label{eq:tpc52-profile-overlap-ratio}
\end{equation}
If \(\Theta(\mathbf W;c)>0\), then
\begin{equation}
 \boxed{
 \frac{\cE_B(\mathbf W;c)^2}
      {\|\cS_{\mathbf W}c\|_2^2}
 =\frac{A_B(\mathbf W)^2}{\Theta(\mathbf W;c)}.}
 \label{eq:tpc52-exact-Theta-ledger}
\end{equation}
If \(c\ne0\) and \(\Theta(\mathbf W;c)=0\), then
\(\cS_{\mathbf W}c=0\) while
\(\cE_B(\mathbf W;c)>0\); the ratio in
\eqref{eq:tpc52-exact-Theta-ledger} is therefore \(+\infty\), in
agreement with the extended-real convention \(A_B^2/0=+\infty\).
For \(c=0\), both sides of the synthesis and envelope identities are
zero, but no condition number and no ratio \(0/0\) is defined.
\end{theorem}

\begin{proof}
Since \(\cI\) is finite, \eqref{eq:tpc52-Mellin-mass-finite} and
\begin{equation}
 \|A_{\bm t}(c)\|_2=\|c\|_2
 \quad\text{for }\mu_{\mathbf W}\text{-almost every }\bm t
 \label{eq:tpc52-unit-layer-norm}
\end{equation}
give Bochner integrability.  In coordinate \(i\), Fubini's theorem and
the product form of \(\mu_{\mathbf W}\) give
\begin{align}
 \int_{\R^q}A_{\bm t}(c)_i\,
      d\mu_{\mathbf W}(\bm t)
 &=\frac{c_i}{(2\pi)^q}
   \int_{\R^q}\prod_{\nu=1}^q
     \widetilde W_\nu(t_\nu)x_{i,\nu}^{it_\nu}
     \,d\bm t\notag\\
 &=c_i\prod_{\nu=1}^qW_\nu(x_{i,\nu}),
 \label{eq:tpc52-coordinatewise-Mellin-inversion}
\end{align}
where the last line is \eqref{eq:tpc52-Mellin-inversion} in each
variable.  This proves \eqref{eq:tpc52-diagonal-output}; orthogonality
of the coordinate basis proves
\eqref{eq:tpc52-diagonal-output-energy}.  Equation
\eqref{eq:tpc52-unit-layer-norm} and the definition
\eqref{eq:tpc52-weighted-Mellin-mass} give
\eqref{eq:tpc52-exact-envelope-equality}.  Finally,
\eqref{eq:tpc52-profile-overlap-ratio} says exactly that
\begin{equation}
 \|\cS_{\mathbf W}c\|_2^2
 =\Theta(\mathbf W;c)\|c\|_2^2.
 \label{eq:tpc52-output-as-overlap}
\end{equation}
Combining this with \eqref{eq:tpc52-exact-envelope-equality} proves
the ledger and all zero-denominator assertions.
\end{proof}

\begin{remark}[What the identity does not say]
\label{rem:tpc52-diagonal-identity-scope}
The coefficient vector \(c\) in
\cref{thm:tpc52-diagonal-Mellin} is the fixed base vector carried by
the physical coordinates.  The theorem does not permit an arbitrary
signed density in \(\bm t\), and it does not assert a lower frame on
the full Mellin coefficient space.  All row, orbit, and residue phases
that have unit modulus may be included in the samples or in \(c_i\);
literal mask zeros may also be included in \(c_i\).  The identity
fails to describe a later map that first sums several physical
coordinates coherently and only then takes a Hilbert norm.
\end{remark}

\subsection{A fixed number of direct residue cells}

The actual provenance uses a fixed finite disjoint union of direct
residue cells.  Their output supports are orthogonal because the orbit
classes are disjoint.  The outer canonical envelope, however, is an
\(\ell^1\)-sum over the cells.  The following comparison records the
resulting fixed constant exactly.

\begin{corollary}[Finite direct-cell \(L^1\) comparison]
\label{cor:tpc52-finite-direct-cell-comparison}
Let \(\mathfrak A\) be a fixed nonempty finite set.  For every
\(a\in\mathfrak A\), let
\(\mathbf W_a=(W_{a,1},\ldots,W_{a,q_a})\) satisfy the hypotheses of
\cref{thm:tpc52-diagonal-Mellin}, let \(c^{(a)}\) be a finite base
vector, and suppose that the corresponding physical outputs lie in
pairwise orthogonal Hilbert subspaces.  Put
\begin{align}
 b_a&:=\|c^{(a)}\|_2,
 &A_a&:=A_B(\mathbf W_a),
 \label{eq:tpc52-direct-cell-base-data}\\
 \cE_B^{\rm dir}
 &:=\sum_{a\in\mathfrak A}A_ab_a,
 &\mathcal B_{\rm dir}^2
 &:=\sum_{a\in\mathfrak A}b_a^2,
 \label{eq:tpc52-direct-cell-envelope}\\
 \Theta_{\rm dir}
 &:=\frac{
   \displaystyle\sum_{a\in\mathfrak A}
       \|\cS_{\mathbf W_a}c^{(a)}\|_2^2}
   {\mathcal B_{\rm dir}^2},
 \label{eq:tpc52-direct-cell-overlap}
\end{align}
whenever \(\mathcal B_{\rm dir}>0\).  Then
\begin{equation}
 A_-^2\mathcal B_{\rm dir}^2
 \leq (\cE_B^{\rm dir})^2
 \leq A_2^2\mathcal B_{\rm dir}^2,
 \qquad
 A_-:=\min_{a\in\mathfrak A}A_a,
 \quad
 A_2^2:=\sum_{a\in\mathfrak A}A_a^2.
 \label{eq:tpc52-direct-cell-envelope-comparison}
\end{equation}
If \(\Theta_{\rm dir}>0\), orthogonality of the cell outputs gives
\begin{equation}
 \boxed{
 \frac{A_-^2}{\Theta_{\rm dir}}
 \leq
 \frac{(\cE_B^{\rm dir})^2}
      {\displaystyle\sum_{a\in\mathfrak A}
       \|\cS_{\mathbf W_a}c^{(a)}\|_2^2}
 \leq
 \frac{A_2^2}{\Theta_{\rm dir}}.}
 \label{eq:tpc52-finite-cell-Theta-ledger}
\end{equation}
If \(\Theta_{\rm dir}=0\), the nonzero base vector has zero physical
output and positive envelope, so the middle ratio is \(+\infty\).
The constants \(A_-\) and \(A_2\) depend only on the fixed direct
cells, shapes, and profile weight, and not on \(X\) or on the physical
coordinate multiplicities.
\end{corollary}

\begin{proof}
Every \(A_a\) is positive by
\eqref{eq:tpc52-Mellin-mass-finite}.  Therefore
\begin{equation}
 \cE_B^{\rm dir}
 =\sum_aA_ab_a
 \geq A_-\sum_ab_a
 \geq A_-\left(\sum_ab_a^2\right)^{1/2}.
 \label{eq:tpc52-direct-cell-lower-L1}
\end{equation}
Cauchy--Schwarz gives
\begin{equation}
 \cE_B^{\rm dir}
 \leq\left(\sum_aA_a^2\right)^{1/2}
      \left(\sum_ab_a^2\right)^{1/2}.
 \label{eq:tpc52-direct-cell-upper-L1}
\end{equation}
These are \eqref{eq:tpc52-direct-cell-envelope-comparison}.  Pairwise
orthogonality and \eqref{eq:tpc52-direct-cell-overlap} give
\begin{equation}
 \sum_a\|\cS_{\mathbf W_a}c^{(a)}\|_2^2
 =\Theta_{\rm dir}\mathcal B_{\rm dir}^2.
 \label{eq:tpc52-direct-cell-output-energy}
\end{equation}
Divide \eqref{eq:tpc52-direct-cell-envelope-comparison} by this
identity.  If \(\Theta_{\rm dir}=0\), every physical cell output is
zero while \(\cE_B^{\rm dir}>0\), proving the last assertion.
\end{proof}

\begin{remark}[Direct cells are essential]
\label{rem:tpc52-direct-cells-essential}
The orthogonality in
\cref{cor:tpc52-finite-direct-cell-comparison} holds for direct residue
classes, such as disjoint classes of \(j\pmod {H_0}\).  It is not
asserted for overlapping signed Fourier cells or for projective-mask
parameters.  Appending cancelling signed layers can change their outer
\(L^1\) envelope without changing the physical multiplier.  Thus the
finite-cell comparison does not license a reverse Minkowski inequality
for an arbitrary signed representation.
\end{remark}

\subsection{Identification with the retained physical coordinates}

For the interface of \cref{sec:tpc52-interface}, take
\begin{equation}
 i=(m,r,\gamma,j),
 \label{eq:tpc52-physical-coordinate-identification}
\end{equation}
include the literal mask, source coefficient, fixed residue amplitude,
and common carrier in \(c_i\), and take the positive samples
\(x_{i,\nu}\) to be the declared arguments of the physical smooth
factors.  Typical examples are \(m_\gamma j/X\) and
\(d_\gamma j/K\).  Before coherent prime summation and before grouping
the residual label, these coordinates are orthogonal by
\eqref{eq:tpc52-exact-physical-energy}.  Hence
\cref{thm:tpc52-diagonal-Mellin,cor:tpc52-finite-direct-cell-comparison}
apply literally to the provenance-fixed direction
\(\cS_X^L\one=G_X^L\).

The only remaining analytic input at this interface is a lower bound
for the explicit nonnegative overlap
\(\Theta_{\rm dir}\).  The diagonal theorem neither assumes nor proves
such a bound.  In particular, it contains no distribution estimate for
primes, no parity-sensitive cancellation, and no fixed-shift
prime-pair conclusion.
